\section{Conclusion}
\label{sec:conclusion}

This paper identified the Selective Deletion Attack---a threat model
in which a system operator exploits individual erasure rights to
distort the statistical history of an AI decision log---and presented
Accountable Redaction, a protocol that makes such distortion
cryptographically detectable.

The core contribution is the $\varepsilon$-statistical consistency
predicate: a ZK-SNARK over Pedersen-committed demographic attributes
that proves, without revealing individual decision content, that no
protected group's conditional approval rate shifted by more than
$\varepsilon$ after a redaction batch.  The protocol separates
authorization ($\pi_{\mathit{auth}}$: was this erasure requested by the
data subject?) from statistical integrity ($\pi_{\mathit{stat}}$: does
this erasure distort the record?)---closing an attack vector absent
from all prior redactable data structures.

Together with the BTV series (Papers~1--2), this work establishes a
three-layer accountability stack for high-stakes AI systems:
no decision without evidence (compile-time),
no delivery without verifiable persistence (runtime),
and no redaction without statistical justification (cryptographic).
The right to be forgotten and the obligation to be accountable are not
irreconcilable---they require only that deletion be provably
non-discriminatory.
